\documentclass{tietreport}

% Import images
\usepackage{graphicx}

% Bibliography configuration
\usepackage[
  date=year,
  style=alphabetic,
  backend=biber,
  sorting=nty,
  maxnames=5,
  minnames=3,
  maxalphanames=4,
  minalphanames=3,
  backref=true,
  doi=false,
  isbn=false,
  url=false,
  eprint=false
]{biblatex}
\addbibresource{references.bib}

%% ==================================================
%% CUSTOMIZE THESE FIELDS FOR YOUR PROJECT
%% ==================================================

% Course/Lab header
\TitlePageHeader{%
  UCS503 - Software Engineering%
}

% Main project title
\ProjectTitle{%
  AgriGuard%
}

% Subtitle or project description
\ProjectSubTitle{%
  AI-Based Crop Disease Detection and Advice for Farmers%
}

% Student information (up to 4 students)
\author{%
  \texttt{1024030114} Vaibhav Chaudhary \\
  \texttt{1024030155} Harshit Kasrija \\
  \texttt{1024030179} Harnoor%
}

% Group number, degree, year, program
\TitlePageSubText{%
  Group: \texttt{XXXX} BE Third Year, CSE%
}

% Faculty advisor/guide name
\AdvisorName{%
  Ms. Paramveer Kaur%
}

% Evaluation period and department info
\TitlePageFooterText{{%
    \large Mid-Semester Evaluation \\[0.2cm]
    \bfseries Computer Science and Engineering
    Department%
  } \\[0.15cm]
  Thapar Institute of Engineering and Technology,
  Patiala%
}

%% ==================================================

\begin{document}

% Generate title page
\maketitle

% Table of contents (auto-generated from chapters)
\tableofcontents

% ===== CHAPTER 1: INTRODUCTION =====

\chapter{Introduction}

\section{Background and Context}

Crop diseases such as leaf blights, rusts and fungal
infections spread quickly and can destroy a large part
of a crop. They cause money loss for farmers and often
lead to too much use of chemicals. At present, most
farmers find out the disease using their own experience,
or they wait for an agriculture expert to visit. This is
slow, and a wrong or late decision can make the damage
worse.

Computer vision and deep learning can help here. Earlier
research has shown that a model trained on leaf photos
can recognise many plant diseases~\cite{mohanty2016}. Our
project, \textbf{AgriGuard}, uses this idea to build a
simple web application. A farmer uploads a photo of an
affected leaf and gets the disease name, how serious it
is, and what to do next. The project also shows how a
machine learning model can be connected with a website,
a database, user login and testing to form one working
system.

\subsection{Problem Statement}

Diagnosis of crop diseases is manual, slow and scattered.
Farmers find it hard to know which disease is present, how
severe it is and which treatment to use. There is a need
for a fast, low-cost and easy tool that can identify the
disease from a photo and explain the next steps in simple
language, without depending on paid services.

\section{Project Objectives}

\begin{enumerate}
  \item \textbf{Disease detection:} Identify the crop
    disease from a leaf photo using a deep learning model
    trained on the PlantVillage dataset, and measure its
    accuracy, precision and recall on a separate test set.
  \item \textbf{Severity and confidence:} Show a confidence
    score and estimate the severity (Low, Medium or High)
    for every result.
  \item \textbf{Advice:} Give simple disease information and
    treatment advice for each detected disease.
  \item \textbf{Records and languages:} Save every report so
    that the farmer can see old results, and support English,
    Hindi and Punjabi in the interface.
  \item \textbf{Free tools:} Build the complete system using
    only free and open-source tools, with a clear separation
    between the website, backend and model.
\end{enumerate}

These objectives are specific and can be checked: the model
results are measured on unseen test images, and the system is
checked with test cases in Chapter~3.

% ===== CHAPTER 2: METHODOLOGY & DESIGN =====

\chapter{Methodology and Design}

\section{System Architecture}

AgriGuard follows a three-tier design: a website for the user,
a backend that does the work, and a database for storage.

\subsection{High-Level Design}

The main parts of the system are:

\begin{itemize}
  \item \textbf{Website (frontend):} Pages for login, scanning
    a crop, viewing results, report history and the AgriBot chat.
  \item \textbf{Backend:} Handles login, checks the uploaded
    image, runs the model, measures severity, prepares the advice
    and saves the report.
  \item \textbf{Database and storage:} Stores users, reports and
    chat messages. Uploaded images are kept as files and linked
    to their reports.
\end{itemize}

For one photo, the system works in these steps:

\begin{enumerate}
  \item The farmer logs in and uploads or captures a leaf photo
    (JPG, PNG or WEBP, up to 8~MB).
  \item The backend checks that the file is a proper image.
  \item The MobileNetV2 model predicts the disease and gives a
    confidence score.
  \item An OpenCV step measures how much of the leaf is damaged
    and gives a severity level.
  \item The result and advice are shown to the farmer and the
    report is saved in the history.
\end{enumerate}

\subsection{Technical Stack}

\begin{itemize}
  \item \textbf{Framework/Language:} Python for the backend
    (FastAPI) and model training, JavaScript with React for the
    website.
  \item \textbf{Database:} SQLite.
  \item \textbf{Tools:} Kaggle notebook with a free GPU for
    training, Git and GitHub for source control, and an AI coding
    assistant (Antigravity) to help build the web application from
    our written specification. The team ran and checked the code.
  \item \textbf{Libraries:} PyTorch~\cite{paszke2019} and
    torchvision for the model, OpenCV and Pillow for images,
    scikit-learn and matplotlib for results, and Tailwind CSS for
    the interface design.
\end{itemize}

All tools are free and open source, and no paid service or
external AI service is needed.

\section{Implementation Details}

\subsection{Module 1: Disease Detection Model}

\paragraph{Dataset.} We used the PlantVillage dataset with color
images~\cite{hughes2015}. Each folder is one class, named after
the crop and the disease, for example
\texttt{Tomato\_\_\_Early\_blight}. The copy we used was already
divided into \texttt{train} and \texttt{val} folders. We kept 15\%
of the \texttt{train} folder (chosen with a fixed random seed) as a
validation set to pick the best training round. The \texttt{val}
folder of the dataset was not used during training and was used
only once at the end as the test set. The application supports 14
crops.

\paragraph{Model.} We used transfer learning. We started with
MobileNetV2~\cite{sandler2018}, which was already trained on
ImageNet, and replaced its last layer so that it gives one output
for each disease class. This model is small and fast, so it could be
trained on a free GPU and can run on an ordinary computer. Training
was done in two phases, as shown in Table~\ref{tab:training}.

\begin{table}[h]
\centering
\begin{tabular}{|p{4.2cm}|p{9cm}|}
  \hline
  \textbf{Setting} & \textbf{Value} \\
  \hline
  Base model & MobileNetV2 (ImageNet weights) \\
  \hline
  Image size, batch size & $224 \times 224$ pixels, 64 \\
  \hline
  Optimizer & AdamW \\
  \hline
  Phase A & 2 epochs, only the new last layer trained
    (learning rate $10^{-3}$) \\
  \hline
  Phase B & 6 epochs, last 6 blocks also trained
    (learning rate $10^{-4}$) \\
  \hline
  Augmentation & Random crop, horizontal and vertical flip,
    rotation up to 20 degrees, color change \\
  \hline
  Best model & Saved at the epoch with the best validation
    accuracy \\
  \hline
\end{tabular}
\caption{Training settings of the disease model}
\label{tab:training}
\end{table}

The best model and the list of class names are saved to files, and
the application loads exactly these files to make predictions.

\subsection{Module 2: Severity Estimation and Treatment Advice}

\paragraph{Severity.} The model gives only the disease name, so
severity is measured separately with OpenCV. The leaf is separated
from the background using color (HSV), the healthy green pixels are
counted, and the remaining leaf pixels are counted as damaged. The
damaged fraction is converted to a level: below 10\% is Low, 10\% to
30\% is Medium, and above 30\% is High. This is a simple
color-based estimate and not a trained model. The application also
shows an overlay on the photo, so the farmer can see which part was
counted.

\paragraph{Advice.} The advice comes from a prepared knowledge base
with one entry for each disease class. It suggests simple and
organic steps first, names chemicals only by their active
ingredient, gives no exact dose, and tells the farmer to read the
product label and ask a local agriculture officer. If the model
confidence is low, the application asks for a clearer photo instead
of giving advice.

\subsection{Module 3: Web Application and User Interface}

The website is made for non-technical users, with large buttons and
a green theme. The login page (Figure~\ref{fig:login}) has one-click
demo accounts for the two roles, Farmer and Expert, along with normal
sign in and account creation. The language can be changed to
English, Hindi or Punjabi from the top bar
(Figure~\ref{fig:lang}).

\begin{figure}[h]
\centering
\includegraphics[width=0.75\textwidth]{images/fig1_login.png}
\caption{Login page with Farmer and Expert demo accounts}
\label{fig:login}
\end{figure}

\begin{figure}[h]
\centering
\includegraphics[width=0.4\textwidth]{images/fig2_language.png}
\caption{Language selector (English, Hindi, Punjabi)}
\label{fig:lang}
\end{figure}

On the Scan Crop page (Figure~\ref{fig:scan}), the farmer can drag
and drop a leaf photo, browse for one, or open the camera. A link
explains which crops are supported and what the model cannot do.
After a photo is chosen, a preview is shown
(Figure~\ref{fig:preview}) and the farmer can change or remove it.

\begin{figure}[h]
\centering
\includegraphics[width=0.85\textwidth]{images/fig3_scan_page.png}
\caption{Scan Crop page for uploading or capturing a leaf photo}
\label{fig:scan}
\end{figure}

\begin{figure}[h]
\centering
\includegraphics[width=0.75\textwidth]{images/fig4_preview.png}
\caption{Photo preview before analysis}
\label{fig:preview}
\end{figure}

Below the upload box, the page shows tips for taking a good photo
and optional boxes for crop, field location and notes
(Figure~\ref{fig:options}). The analysis starts when the farmer
presses \textit{Analyze Leaf Now}.

\begin{figure}[h]
\centering
\includegraphics[width=0.8\textwidth]{images/fig5_options.png}
\caption{Photo tips, optional details and the Analyze Leaf Now button}
\label{fig:options}
\end{figure}

The result page (Figure~\ref{fig:result}) shows the crop, the type
of disease, the disease name with its scientific name, a short
description, the severity with the percentage of damaged area, and the
confidence of the model. An overlay marks the damaged parts of the
leaf, and the farmer can switch back to the original photo. A button
opens AgriBot for the same diagnosis.

\begin{figure}[h]
\centering
\includegraphics[width=0.9\textwidth]{images/fig6_result.png}
\caption{Result page for a tomato leaf: Early blight, High severity
(34.1\% damaged area), 98\% confidence}
\label{fig:result}
\end{figure}

\subsection{Module 4: AgriBot Chatbot (Current Status)}

AgriBot is meant to answer the farmer's questions about the detected
disease and the action needed. The chat page, the button on the
result page and the saving of chat history are in place. However, the
chatbot is still a prototype: we have not trained or tested it
properly yet, so its answers can be short or incomplete, and we have
not measured how good it is. For this reason AgriBot is not included
in the evaluation in Chapter~3. Improving it is the first item of
future work.

% ===== CHAPTER 3: RESULTS & EVALUATION =====

\chapter{Results and Evaluation}

\section{Testing Methodology}

We checked the system in the following ways:

\begin{itemize}
  \item \textbf{Model testing:} The trained model was tested on
    images from the held-out test set that it never saw during
    training. We measured accuracy, precision, recall and F1 score
    for each class.
  \item \textbf{Functional testing:} We tried the main steps of the
    application by hand: login, uploading photos, viewing the
    result, report history and language change.
  \item \textbf{Input checking:} We tried a blurry photo, a photo
    that is not a leaf, a wrong file type and a file larger than
    8~MB, to see that the application gives a proper message.
  \item \textbf{Performance testing:} We noted the time the
    application takes to show a result after a photo is uploaded.
\end{itemize}

\section{Performance Metrics}

Table~\ref{tab:model} shows the results of the model on the test
set. The values are taken from the report saved after training.

\begin{table}[h]
\centering
\begin{tabular}{|p{6.5cm}|p{6.5cm}|}
  \hline
  \textbf{Item} & \textbf{Result} \\
  \hline
  Number of classes & \textbf{[FILL]} \\
  \hline
  Train / validation / test images & \textbf{[FILL]} \\
  \hline
  Test accuracy & \textbf{[FILL]} \\
  \hline
  Precision (macro average) & \textbf{[FILL]} \\
  \hline
  Recall (macro average) & \textbf{[FILL]} \\
  \hline
  F1 score (macro average) & \textbf{[FILL]} \\
  \hline
\end{tabular}
\caption{Model results on the held-out test set}
\label{tab:model}
\end{table}

Table~\ref{tab:example} shows an example result from the running
application (see also Figure~\ref{fig:result}).

\begin{table}[h]
\centering
\begin{tabular}{|p{4.5cm}|p{8.5cm}|}
  \hline
  \textbf{Item} & \textbf{Output} \\
  \hline
  Predicted disease & Tomato, Early blight
    (\textit{Alternaria solani}), a fungal disease \\
  \hline
  Confidence & 98\% \\
  \hline
  Severity & High (34.1\% of the leaf area damaged) \\
  \hline
\end{tabular}
\caption{Example output of the application}
\label{tab:example}
\end{table}

Table~\ref{tab:tests} lists the functional tests we carried out.

\begin{table}[h]
\centering
\begin{tabular}{|p{4.2cm}|p{6.3cm}|p{2.2cm}|}
  \hline
  \textbf{Test} & \textbf{What should happen} &
  \textbf{Result} \\
  \hline
  Good leaf photo & Disease, confidence and severity are shown and
    the report is saved & \textbf{[FILL]} \\
  \hline
  Blurry or non-leaf photo & The application asks for a clearer
    photo & \textbf{[FILL]} \\
  \hline
  Wrong file type or file over 8~MB & The photo is rejected with a
    message & \textbf{[FILL]} \\
  \hline
  Login as Farmer and Expert & Correct pages for each role &
    \textbf{[FILL]} \\
  \hline
  Old reports & Past reports are listed after login &
    \textbf{[FILL]} \\
  \hline
  Language change & The interface changes to Hindi and Punjabi &
    \textbf{[FILL]} \\
  \hline
  Response time & Result shown within a few seconds &
    \textbf{[FILL]} \\
  \hline
\end{tabular}
\caption{Functional tests of the application}
\label{tab:tests}
\end{table}

\section{Analysis}

\begin{itemize}
  \item \textbf{Did we meet our objectives?} The main workflow
    works from start to end: a farmer uploads a leaf photo and gets
    the disease, confidence, severity and advice, and the report is
    saved. The interface supports three languages. The chatbot part
    of the advice objective is not finished yet.
  \item \textbf{Where did we do well?} A small pre-trained model
    gave good results with a short training time on a free GPU, and
    the whole system runs on free tools.
  \item \textbf{What were the challenges?} The PlantVillage photos
    are taken on plain backgrounds and many images show the same
    leaf, so the test accuracy is likely higher than what we would
    get on real farm photos. The severity is based on color, so
    light and background can change it. The model only knows the
    crops it was trained on, so a leaf from another crop can get a
    wrong answer.
  \item \textbf{How did we handle them?} The application shows
    the supported crops, asks for a clearer photo when the
    confidence is low, and shows a warning that an agriculture
    expert should check the treatment advice before it is used on
    a real farm.
\end{itemize}

% ===== CHAPTER 4: CONCLUSION =====

\chapter{Conclusion}

\section{Summary of Work}

We built AgriGuard, a web application that finds crop diseases
from leaf photos. We trained a MobileNetV2 model on the PlantVillage
dataset, added an OpenCV step to estimate severity, prepared a
knowledge base for advice, and built a multilingual website with
login, report history and a chat page for AgriBot. Most of our
objectives were met. The AgriBot chatbot is still a prototype.

\section{Key Findings}

\begin{itemize}
  \item \textbf{Finding 1:} Transfer learning with a small model
    gives good disease recognition on PlantVillage with a short
    training time (see Table~\ref{tab:model}).
  \item \textbf{Finding 2:} A disease model alone is not enough for
    farmers. Severity, confidence, clear advice and saved reports
    make the result useful.
  \item \textbf{Finding 3:} Results on a clean dataset do not
    guarantee the same results on real farm photos, so testing on
    real photos is important.
\end{itemize}

\section{Future Work and Recommendations}

\begin{enumerate}
  \item Finish and test AgriBot properly, using a checked
    knowledge base.
  \item Train with real farm photos, add more crops, and reject
    photos that are not leaves.
  \item Use a trained model for severity instead of the color
    method.
  \item Complete the expert review feature, where agriculture
    officers can verify reports.
  \item Make a mobile application and put the system online with a
    permanent database.
\end{enumerate}

\section{Final Remarks}

This project taught us how to take a machine learning model from
training to a working application, and how much the surrounding
parts (image checks, user roles, storage and testing) matter. With
checked advice and real farm data, a tool like AgriGuard could
become useful for farmers.

% ===== BIBLIOGRAPHY =====

% Print bibliography with heading in TOC
\printbibliography[heading=bibintoc]

\end{document}
